\section{Short paths and uniform insertion estimates}\label{sec:paths}

Proposition~\ref{prop:local-energy} determines the occupation outside the
$R$-cell constants.  We now prepare comparisons that will control the
remaining variation between cells.  For a cell $v$ of $\cC_d$ or
$\cC_R$, define a measure on the remaining $N-1$ coordinates $Y'$ by
\begin{equation}\label{eq:palm-density}
 P_v(\dd Y')=s_v(Y')^2\dd Y',\qquad
 s_v(Y')=\left(\frac V{|v|}\int_v\Phi(y,Y')^2\,\dd y\right)^{1/2}.
\end{equation}
Translation invariance makes the first particle uniform on the torus.
Thus $P_v$ is the law obtained by conditioning that particle to lie in
$v$ and then forgetting its position.  We call it the deletion measure
for $v$.

We shall prove two comparisons of these square-root densities.  For the
uniform average over ordered adjacent $d$-cell pairs, the local target is
\begin{equation}\label{eq:local-fourth}
 \E_{v\sim w'}\int
        \frac{|s_v-s_{w'}|^4}{s_v^2+s_{w'}^2}\,\dd Y'
                         \le C_jD^{-1-u/8},
\end{equation}
where the quotient is zero when both densities vanish.  For the distant
target, let $\mathcal P_R$ be the ordered pairs $(v,w')\in\cC_R^2$
whose centers have forward first-coordinate displacement in
$[.30K,.45K]$, with no restriction on the transverse displacements.
We shall prove
\begin{equation}\label{eq:distant-comparison}
 \E_{(v,w')\in\mathcal P_R}\int |s_v-s_{w'}|^2\,\dd Y'
       =o(D^{-1})+o_K(1).
\end{equation}
The first estimate controls the passage from fine to coarse root mean
squares; the second controls their variation across the full torus.
Section~\ref{sec:assembly} transfers both estimates to the arithmetic
means defining the cell projections.

The constructions in Sections~\ref{sec:local-palm} and
\ref{sec:long-palm} use two copies of the ground-state paths.  They
change trajectories while retaining their endpoints and arrange that
deleting the two distinguished midpoints leaves the same remaining
coordinates.  This section constructs the exact conditional path law and
establishes insertion estimates uniform in the other allowed paths and
in the removal of interaction checks.  Those uniformities will permit
the later likelihood comparisons.  The construction adapts the
short-trajectory method of~\cite[Section~5]{D}; for the present bounded
interaction, its exact conditional weight must be retained.

Fix $x>0$ sufficiently small after $u,w$, and put
\begin{equation}\label{eq:path-scales}
 t=\frac1{x\log D},\qquad f_0=D^{-x},\qquad s_1=D^{-4/5},
 \qquad b_0=\frac{C_1}{\sqrt x},\qquad L_0=C_2\sqrt{\log D}.
\end{equation}
Choose the large absolute constants in the order $C_1$, then $C_2$.
We may decrease $x$ so that $b_0$ exceeds every absolute step size used
below.  Constants denoted $C_x$ may depend on $x$ and on fixed
neighborhood sizes.  In exponents such as $Cx$ and $C\sqrt x$, however,
$C$ is independent of small $x$.  All assertions hold for sufficiently
large $D$ and then sufficiently large $K$, in the order of limits fixed
in Section~\ref{sec:setup}.

\subsection{Exposed data and the exact conditional law}

Restrict the stationary ground-path law of Section~\ref{sec:setup} to
$[-t,t]\subset[-1/2,1/2]$.  Write
\[
 Y_i=\omega_i(-t),\qquad X_i=\omega_i(0),\qquad Z_i=\omega_i(t).
\]
A lift $\widetilde\omega_i$ to $\R^3$ satisfies the motion cutoff if
\begin{equation}\label{eq:motion-cutoff}
 |\widetilde\omega_i(q)-\widetilde\omega_i(q')|
 \le L_0\sqrt{|q-q'|+r^2},\qquad -t\le q,q'\le t.
\end{equation}
Call a label \emph{eligible} when the following conditions hold:
\begin{enumerate}
\item At each endpoint time its position is more than $s_1$ from every
other label's position.  The two endpoints have representatives at
Euclidean distance at most $b_0$.
\item The path has the winding determined by these nearby representatives,
satisfies \eqref{eq:motion-cutoff}, and remains within $20b_0$ of their
mean $m_i=(Y_i+Z_i)/2$ in this lift.
\end{enumerate}
For sufficiently large $K$, these nearby displacements and all subsequent
local comparisons have unique compatible lifts, up to a common torus
translation.

Independently flag each label with probability $f_0$, and let $I$ be the
set of flagged eligible labels.  We expose the two endpoint
configurations, all flags, membership in $I$, and the complete paths of
labels outside $I$.  We call those paths \emph{frozen}.  A further
augmentation turns the soft interaction into an exact indicator
constraint.  Set
\[
 a_{ik}(\omega)=\exp\left\{-\int_{-t}^t
                   W(\omega_i(q)-\omega_k(q))\,\dd q\right\}.
\]
In the Feynman--Kac density, replace each factor $a_{ik}(\omega)$ by
integration over $u_{ik}\in[0,1]$ of the indicator
\begin{equation}\label{eq:soft-check}
                    \1_{\{u_{ik}\le a_{ik}(\omega)\}}.
\end{equation}
Thus the joint density is the original endpoint-weighted free-path
density times these indicators and Lebesgue measure in all $u_{ik}$.
Its path marginal is unchanged.  Equivalently, given the paths, one may
sample $u_{ik}$ uniformly on $[0,a_{ik}(\omega)]$.  These variables are
not independent uniforms on $[0,1]$ after conditioning on the paths.
Expose all of them, and denote the entire exposed datum by $\cD$.

Let $B_i$ be the probability law of the Euclidean Brownian bridge of
generator $\Delta$ on $[-t,t]$ between the nearby representatives of
$Y_i,Z_i$.  Its path space uses the chosen lift.  For $i\in I$, let $J_i$
be the indicator of the motion cutoff and the $20b_0$ tube.

\begin{lemma}[Conditional soft-interaction law]\label{lem:conditional-path-law}
For almost every exposed datum $\cD$, the conditional law of the paths
with labels in $I$ is
\begin{equation}\label{eq:conditional-path-law}
 \frac1{Z(\cD)}\prod_{i\in I}J_i(\omega_i)
 \prod_{\substack{i<k\\\{i,k\}\cap I\ne\varnothing}}
       \1_{\{u_{ik}\le a_{ik}(\omega)\}}
                    \prod_{i\in I}B_i(\dd\omega_i),
 \qquad Z(\cD)>0.
\end{equation}
Frozen paths in this formula are held at their exposed values.  Every
pair indicator can fail only if the two paths approach within distance
$r$ at some time.
\end{lemma}

\begin{proof}
At fixed endpoints, the factors involving the ground function are
constant.  Conditional on winding, each free torus bridge is a
Euclidean bridge; its winding weight also becomes constant.  Endpoint
separation is already determined by the exposed endpoints.  Given these
endpoints, the eligibility of any one label depends only on that label's
path.  For a flagged label outside $I$ its path, and hence the failure of
eligibility, is exposed.  For a label in $I$ eligibility imposes exactly
$J_i$ and the selected winding.  Disintegration of the augmented density
therefore gives \eqref{eq:conditional-path-law}.  Its normalizer is
positive on almost every realized fiber.  Finally, $W$ vanishes outside
range $r$, so absence of a close approach makes $a_{ik}=1$.
\end{proof}

All subsequent conditional comparisons use this indicator law.  The
values $u_{ik}$ remain fixed throughout a comparison; collision bounds
will hold uniformly in them.

\subsection{Good sites and their probability}

For an interior choice at a lattice site $z\in\cL$, we use labels in $I$
whose initial endpoint belongs to $U_z$ and whose endpoint displacement
is at most one.  Later, a label at the end of a comparison may have
displacement up to $b_0$; its site will be the unit cell containing $m_i$.

Choose a uniform mesh on each half of $[-t,t]$, with an even number of
intervals and common interval length $\Delta\in[r^2/2,r^2]$.  The
midpoints of the two halves are therefore mesh points as well.  For a
mesh interval $e$, define
\[
                   A_e(\omega)=r^{-1}\operatorname{diam}
                                             \widetilde\omega(e).
\]
Fix neighborhoods large enough to contain every collision test for a
label associated with $z$.  Balls of radius $1000b_0$ about $z$, with a
fixed enlargement to include intersecting unit cells, suffice.  The
actual collision tests below use distances less than $200b_0$.
Fixed further enlargements do not affect any conclusion.
The flagged paths are sparse, so endpoint counts will suffice to sum
their collision costs.  For the denser frozen collection we also record
positions and short-interval oscillations, allowing a sharper sum.

A site $z$ is \emph{good} if the following conditions hold:
\begin{enumerate}
\item It has at least $cf_0D$ interior labels.
\item In its neighborhood, every unit cell has at most $CD$ endpoints
at each endpoint time and at most $Cf_0D$ flagged endpoints.  At each
endpoint time, every ball of radius $h\in[r,1]$ in the neighborhood has
flag count at most
\begin{equation}\label{eq:flag-count}
                         C(f_0Dh^3+D^{.01}).
\end{equation}
\item Each frozen path visiting the neighborhood satisfies
\eqref{eq:motion-cutoff} and has lifted diameter at most $10b_0$.
For every mesh point $q$, either adjacent interval $e$, and every ball
of radius $h\in[r,1]$ there, the sum over frozen paths whose positions at
$q$ are in the ball, with weights $(1+A_e)^3$, is at most
\begin{equation}\label{eq:weighted-count}
                            C(Dh^3+D^{.01}).
\end{equation}
\end{enumerate}
Here $c$ is sufficiently small and $C$ sufficiently large.  These are
conditions on $\cD$.  To make their measurability explicit, one may use
countably many ball centers and radii.  Equivalently, strengthened
finite tests on enlarged ball nets at dyadic scales imply the stated
bounds with a fixed change in $C$.

We also define an event concerning the \emph{original sampled paths} in
a fixed neighborhood: require \eqref{eq:weighted-count} for all paths,
require \eqref{eq:motion-cutoff} for every visiting path, and require its
supremum distance from each of its endpoints to be at most $b_0$ in the
lift.  We call these the \emph{full-path bounds}.  They will estimate the
mass of successful comparisons.  We do not add this event to $\cD$ or
condition the bridge law on it.

\begin{proposition}[Probability of good data]\label{prop:good-data}
The constants in the preceding definitions can be chosen so that
\begin{equation}\label{eq:good-prob}
\begin{aligned}
 \Pp(z\text{ is bad})&\le C_xD^{-1-1/8},\\
 \Pp(\text{every site in }A\text{ is bad})&\le p^{|A|}
                 &&(A\subset\cL),\\
 \Pp(\text{full-path bounds fail in a fixed neighborhood})
                  &\le C_xD^{-6}.
\end{aligned}
\end{equation}
In the second line $p>0$ may be any sufficiently small fixed number,
after increasing the threshold for $D$.  The constants and thresholds
are uniform in the position of the neighborhood and in the volume.
\end{proposition}

\begin{proof}
We give separate estimates for motion, the supply of eligible labels,
and the fine spatial counts.  All use the deletion estimate
\eqref{eq:path-domination} and the count exponential bound
\eqref{eq:count-mgf}.

\paragraph{Motion and exceptional visitors.}
Gaussian maximal bounds imply that a free path on the short slab fails
\eqref{eq:motion-cutoff}, or has supremum distance greater than $b_0$
from either endpoint, with probability at most $CD^{-40}$.  For the
increment condition, apply Gaussian bounds with fixed slack to all mesh
point pairs, and bound the oscillation on every intervening interval.
There are polynomially many such tests.  The choice of $C_1$, followed
by $C_2$, makes the exponent as large as required.

The same bound applies to the individual cutoffs of a free trial bridge
with endpoint displacement at most $b_0$, and also to a bridge pinned
at time zero at a point within $8b_0$ of its endpoint mean.  For this
last assertion the cutoffs are \eqref{eq:motion-cutoff} and the
$20b_0$ tube, not the stronger distance-$b_0$ condition.  The prescribed
linear segments have increment coefficient at most
$Cb_0/\sqrt t=CC_1\sqrt{\log D}$ and lie well inside that tube.
A centered bridge is Brownian motion minus its linear interpolation,
so Gaussian maximal bounds apply after enlarging constants.  Splitting
an interval that crosses the pin gives the same estimate there.

Here is the transfer from free exceptions to ground-law visitors.  For
a site set $A$, let $N_A(\omega)$ count how many of the fixed
neighborhoods centered in $A$ a path visits.  Suppose a
reversal-invariant single-path event $\mathcal E$ has free probability
at most $l$, uniformly in the starting point.  For each fixed $q>0$,
\[
 \E^y_{\rm free}\left[e^{qN_A\1_{\mathcal E}}-1\right]
          \le C_q l^{1/3}e^{-d_K(y,A)}.
\]
To justify the uniform exponential moment of $N_A$, mark successive
exits from unit-radius balls centered at the previous exit point.
Each segment meets only a bounded number of the prescribed
neighborhoods.  If $\tau$ is one exit time, the strong Markov property
and exponential Markov inequality bound the probability of at least
$m$ exits in time at most two by
$e^{2h}(\E e^{-h\tau})^m$.  Taking $h$ large yields every fixed
exponential moment.  Reaching the neighborhoods of $A$ also has a
Gaussian distance tail.  H\"older's inequality combines these two
facts with $\Pp(\mathcal E)\le l$ to prove the displayed estimate.

Expand the product over labels of
$1+(e^{qN_A\1_{\mathcal E}}-1)$, apply
\eqref{eq:path-domination} to each selected set, and then use
\eqref{eq:count-mgf}.  For sufficiently small $l$ this gives
\begin{equation}\label{eq:visitor-exponential}
 \log\E\exp\left(q\sum_iN_A(\omega_i)\1_{\mathcal E_i}\right)
                    \le C_qDl^{1/3}|A|.
\end{equation}
The same argument applies to an event after adjoining its time
reversal.  It also bounds the expected number of exceptional visitors:
$CDl^{1/3}$ per site.  With $l\le CD^{-40}$, these losses are smaller
than the powers required in \eqref{eq:good-prob}.  Choosing $q$ first,
then $D$, makes the joint exponential base as small as desired.
The endpoint displacement exceeding one has free probability tending
to zero.  Applying \eqref{eq:visitor-exponential} to that event shows
that loss of a fixed positive fraction of $D$ labels at every site of
a prescribed set costs $e^{-cD|A|}$.

\paragraph{Eligible labels and their flags.}
The lower count tail in \eqref{eq:count-tails} supplies a fixed positive
multiple of $D$ initial endpoints at each site, outside its stated
exceptional event.  The preceding motion estimates remove only a
small fraction.  By \eqref{eq:pair-bound}, the expected number of
ordered close endpoint pairs in a fixed neighborhood is at most
$C_xD^2s_1^3$.  Therefore the probability that endpoint separation
removes a fixed fraction of the supply is at most
\[
                         C_xDs_1^3=C_xD^{-7/5}.
\]
Use this at each of the two endpoint times, with the fixed enlargement
needed after displacement removals.

For a joint bound, it is enough to bound the event of at least one
close pair in each of $m$ sufficiently separated neighborhoods at one common endpoint time.
Its probability is at most
\begin{equation}\label{eq:joint-close-pairs}
                            (C_xD^2s_1^3)^m.
\end{equation}
Indeed, at a common time the spatially separated neighborhoods
ensure that the chosen pairs have distinct labels across neighborhoods.  Repeat the cube-cover proof of
\eqref{eq:pair-bound}, using \eqref{eq:path-domination} for the $2m$
labels.  The remaining start weights are bounded by
$\E\prod_{z\in A}F_z^2\le(CD)^{2m}$, where
$F_z=\sum_i e^{-d_K(Y_i,z)}$.  For completeness,
\[
 \sum_{z\in A} e^{-d_K(y,z)}
                         \le Ce^{-d_K(y,A)/2}.
\]
The proof of \eqref{eq:count-mgf} also holds with distance exponent
$1/2$, so
$\log\E e^{h\sum_{z\in A}F_z}\le ChDm$ for fixed sufficiently
small $h$.  The elementary bound
$\prod F_z^2\le(\sum F_z/m)^{2m}$, followed by the exponential
moment bound at $h$ of order $D^{-1}$, gives the asserted product
moment.  This proves \eqref{eq:joint-close-pairs}; its base tends to
zero since $D^2s_1^3=D^{-2/5}$.

A positive fraction of $D$ eligible labels consequently remains with
the required one-site and joint bounds.  Conditional on all paths,
independent flags select at least $cf_0D$ of them except with a
binomial lower-tail cost.  Lists based at different sites are disjoint,
so the same binomial argument gives joint bounds.  Upper unit counts
follow from \eqref{eq:count-tails}; their flagged versions follow from
binomial upper tails after the unit counts have been bounded.

\paragraph{Fine spatial counts.}
It remains to establish \eqref{eq:flag-count} and
\eqref{eq:weighted-count}, including the latter for all original
paths.  At sufficiently separated sites choose one test per site:
a radius $h_z\in[r,1]$, a mesh time, and, for the weighted test, an
adjacent interval.  In each count retain only paths satisfying the
motion and diameter cutoffs.  Exceptional visitors have already been
estimated.  Separation ensures that a retained path contributes to
at most one chosen site's test.  On the cutoff,
$A_e\le C\sqrt{\log D}$.  With $b_D'=(\log D)^{-2}$, a one-path
exponential increment is therefore bounded by a sum of
\[
 Cb_D'(1+A_e)^3
             \1_{\{\omega(q_z)\text{ lies in the tested ball}\}}.
\]
Consider its free expectation on the longer slab $[-1,1]$.  The
observation time remains uniformly away from either endpoint.  The
heat kernel and Gaussian increment moments give, for each site,
\[
                       Cb_D'h_z^3 e^{-3d_K(y,z)}.
\]
For a preceding interval, first use the heat kernel at its initial
point, translate the ball by the increment, and integrate the
increment with its oscillation weight.  This proves the same bound;
the time-reversed tests satisfy it as well.  For a flagged endpoint
count, averaging the independent flag supplies the additional factor
$f_0$.

Expand over labels and apply \eqref{eq:path-domination} as before.
Equation~\eqref{eq:count-mgf} then gives logarithmic moment bound
$Cb_D'D\sum_z h_z^3$, or $Cb_D'f_0D\sum_z h_z^3$ for flag counts.
For unequal radii, write $h_z^3$ as an integral of level-set indicators
on $[0,1]$, use convexity, and use
$\sum_{z\in A}e^{-3d_K(y,z)}\le Ce^{-d_K(y,A)}$ for each level
set.  Exponential Markov inequality, with the constant in the tests
chosen sufficiently large, bounds simultaneous violations by
\[
                          e^{-cb_D'D^{.01}|A|}.
\]
There are only polynomially many tests per site on enlarged dyadic
ball nets; this cost still tends to zero faster than every negative
power of $D$.

Finally, partition arbitrary site sets into a bounded number of
separated classes.  Split failures among the finitely many types just
considered.  For any assignment of types and classes, one class and
type contains a fixed fraction of the failing sites, and its
probability has arbitrarily small exponential base.  The finite
number of assignments is absorbed by making that base smaller.
The one-site error is dominated by $D^{-7/5}$, and the full-path
errors have the stronger bounds proved above.  This yields all of
\eqref{eq:good-prob}.
\end{proof}

\subsection{The cost of one pair check}

We now work at fixed good data.  Paths of other labels in $I$ may take
arbitrary values satisfying their individual cutoffs; they need not be
sampled from the conditional law.  This uniformity will allow us to
change subsets of pair checks in later comparisons.

For distinct labels $i,k$, with trial label $i$ eligible, define
\begin{equation}\label{eq:beta}
\begin{split}
 \beta_{ik}&=b_D\left(1+
    \frac1{s_1\vee d_K(Y_i,Y_k)}+
    \frac1{s_1\vee d_K(Z_i,Z_k)}\right),\\
 b_D&=C_xD^{-1}(\log D)^C.
\end{split}
\end{equation}
The fixed polylogarithmic factor may be enlarged throughout this
section.  In particular, $\beta_{ik}\le C_xD^{-.19}$ for large $D$.

\begin{lemma}[Collision costs]\label{lem:collision-cost}
Fix a trial label $i$ with eligible endpoints and a distinct opponent
$k$ satisfying the motion cutoff.  The $B_i$-probability that the trial
satisfies its individual cutoffs and comes within $r$ of the opponent
is at most $C\beta_{ik}$.  The same holds for a trial pinned at a point
$y$ within $8b_0$ of $m_i$, provided
$d_K(y,\omega_k(0))>s_1$, with the additional cost
\begin{equation}\label{eq:midpoint-cost}
                   m(y,\omega_k)=
                       \frac{b_D}{s_1\vee d_K(y,\omega_k(0))}.
\end{equation}
If instead $y$ is uniform in a unit cell at an absolute bounded
distance from $m_i$, entirely within this pin range, the integrated
cost is at most $C_x\beta_{ik}$, including failure of the midpoint
separation condition.  These bounds also bound failures of
\eqref{eq:soft-check}, uniformly in its exposed auxiliary variable.
\end{lemma}

\begin{proof}
Orient each half of an unpinned bridge from its nearer endpoint.  For
a pinned bridge, do the same on the two halves of each pinned leg.
At the starting point of an oriented piece, suppose the separation
from the opponent is $h\ge s_1$.  On the motion cutoffs the paths
cannot collide before elapsed time
$c_xh^2/\log D$.  This is much greater than $\Delta$.  A collision on
a mesh interval forces the positions at an adjacent mesh point to be
within $Cr\sqrt{\log D}$.  At elapsed time $q$ in this oriented
piece the trial has density at most $Cq^{-3/2}$, since we use at most
half of the relevant bridge.  Summing from the lower cutoff gives
\[
 \begin{split}
 &C_xr^3(\log D)^{3/2}
       \sum_{q\ge c_xh^2/\log D}q^{-3/2}\\
 &\hspace{25mm}\le
       C_x\frac{r(\log D)^C}{h},
 \end{split}
\]
where the sum is over mesh points and $\Delta\asymp r^2$.
All tests involve nearby lifts; alternatively the periodized Gaussian
bound gives the same estimate.  Adding the oriented pieces proves
\eqref{eq:beta} and \eqref{eq:midpoint-cost}.  Uniform averaging over a
unit cell integrates the inverse-distance singularity, and midpoint
separation fails on volume at most $Cs_1^3\le Cb_D$.  This proves the
averaged assertion.  The last assertion follows from
Lemma~\ref{lem:conditional-path-law}.
\end{proof}

\subsection{Uniform insertion probabilities}

The one-check bound is already enough for the sparse set of flagged
opponents.  The frozen paths are much more numerous.  Their weighted
mesh counts give a sharper sum, of order $rDt$, that still tends to
zero on the chosen short slab.

\begin{proposition}[Insertion at good data]\label{prop:insertion}
Fix exposed data for which the sites involved are good, with all
required neighborhoods included in their tests.  The following bounds
hold uniformly after deletion of any subset of pair checks.
\begin{enumerate}
\item A $B_i$ trial with label $i\in I$ satisfies its individual
cutoffs and all imposed checks against frozen paths and specified
other labels of $I$ with probability at least $3/4$.  The other
unexposed paths may have any values satisfying their individual
cutoffs.
\item The same bound holds when one first chooses a uniform midpoint
in a unit cell at an absolute bounded distance from the trial mean,
entirely within pin range, and then chooses the two independent
bridge legs.  Midpoint separation by $s_1$ from every opponent may
be imposed as an additional test.  For two such trials in different
copies, using the same uniform midpoint and conditionally independent
legs, joint success in their respective copies has probability at
least $1/2$.
\item Suppose the opponents are paths from the original sample
satisfying the full-path bounds in the relevant neighborhood.  A
trial pinned at any point in its pin range, separated by $s_1$ from
all opponent midpoints, has success probability at least $3/4$,
provided its endpoints have the stated separations.  This remains
true after removing opponents or adding or replacing boundedly many
relevant paths satisfying the individual cutoffs and these
separations.
\end{enumerate}
No assertion requires checks among the specified opponent paths to
hold.  Only their individual cutoffs and the stated data bounds are
used.
\end{proposition}

\begin{proof}
For opponents in $I$, tubes restrict attention to endpoints in a fixed
neighborhood.  Dyadic summation of \eqref{eq:flag-count} gives
\[
 \sum_k\beta_{ik}
       \le C_xb_D\left(f_0D+\frac{D^{.01}}{s_1}\right)=o(1).
\]
Lemma~\ref{lem:collision-cost} proves this for both the free and
uniform-midpoint trials, including the latter's separation test.

For frozen paths, work on one oriented piece of a trial as in the
preceding lemma.  Its starting point is $s_1$-separated from all
opponents; for a pinned piece this includes midpoint separation.
The motion cutoffs exclude collisions until elapsed time
\[
                        q_0=c_xs_1^2/\log D.
\]
At the initial mesh point of an interval with elapsed time $q\ge q_0$,
the trial density is bounded in its lift by
\[
                  Cq^{-3/2}e^{-c|z-a_q|^2/q},
\]
where $a_q$ is the bridge mean.  Conditional on that position, its
motion on the interval has a deterministic drift toward the other
bridge endpoint plus an independent centered Gaussian part.  On the
trial confinement event the drift is bounded by
$C_x\Delta/t\le r$ for large $D$.  The supremum of the centered
part divided by $r$, denoted $U$, has a Gaussian tail independent of
the conditioned position.  A collision with frozen path $k$ thus
requires the trial's mesh-point position to be within
\[
                        Cr(1+A_e(\omega_k)+U)
\]
of that opponent's mesh-point position.

We may now drop the trial cutoffs when integrating this necessary
event.  Truncate $U$ at a sufficiently large multiple of
$\sqrt{\log D}$, incurring a negligible error even after summing all
intervals.  The displayed radius is $o(\sqrt q)$, uniformly for
$q\ge q_0$.  Integrating the Gaussian density and $U$ and summing
over opponents gives a cost per interval at most
\begin{equation}\label{eq:frozen-one-interval}
\begin{split}
 &Cr^3q^{-3/2}\sum_k(1+A_e(\omega_k))^3
       e^{-c'd_K(\omega_k(q_{\rm mesh}),a_q)^2/q}\\
 &\hspace{30mm}\le C_xr^3\bigl(D+D^{.01}q^{-3/2}\bigr).
\end{split}
\end{equation}
Only paths visiting the tested neighborhood occur.  To obtain the
last inequality, use \eqref{eq:weighted-count} in concentric shells
of radius $\sqrt q$, and cover larger shells by unit balls if needed.
Since $q_0\gg r^2$, all the smallest relevant radii are among the
tested scales.  Gaussian decay makes the shell sum convergent.

Summing \eqref{eq:frozen-one-interval} over the $O(t/\Delta)$ mesh
intervals on all oriented pieces yields
\begin{equation}\label{eq:frozen-insertion-cost}
 C_xr\left(Dt+\frac{D^{.01}}{\sqrt{q_0}}\right)+o(1)=o(1).
\end{equation}
Indeed $rDt=O(t)\to0$ and
$rD^{.01}/\sqrt{q_0}=O_x(D^{-.19}\sqrt{\log D})$.
For a uniform midpoint, its separation failure from frozen paths
costs at most $C_xDs_1^3=o(1)$, because their number in the relevant
neighborhood is $O_x(D)$ by \eqref{eq:weighted-count}.

The trial's individual cutoffs fail with probability $O(D^{-40})$,
as proved above.  Thus the total failure probability in Assertions~1
and~2 is $o(1)$ and in particular at most $1/4$.  The union bound for
two copies gives joint success at least $1/2$, even though their
midpoints are shared.

For Assertion~3, the full-path bounds allow
\eqref{eq:frozen-one-interval} to be applied to all original opponents.
Removal only decreases the cost.  Each additional relevant path costs
at most $C_xb_D/s_1=o(1)$ by
Lemma~\ref{lem:collision-cost}; their number is bounded independently
of $D,K$.  The same success bound follows.  Every estimate bounds a
sum of potentially failing individual or pair tests, so deleting
checks and allowing previously untested checks among opponents does
not change the proof.
\end{proof}

\begin{corollary}[Domination by free bridges]\label{cor:bridge-domination}
Take finitely many labels at good sites, with product reference law
$\prod_iB_i$, individual cutoffs, and any selection of the pair checks
just considered.  Normalize their indicator-weighted law.  Its
marginal on any $q$ labels has density at most $C^q$ relative to their
product free-bridge law.  The same conclusion holds conditionally on
any individually allowed values of the other paths for which the
conditional law is defined.
\end{corollary}

\begin{proof}
For one label, the conditional density is an indicator divided by
its insertion probability, which is at least $3/4$ by
Proposition~\ref{prop:insertion}.  Delete the $q$ selected labels and
then insert them successively.  Each insertion changes the
normalization comparison by at most $4/3$.  The resulting measure
domination, tested against an arbitrary nonnegative function of those
$q$ labels, proves the assertion.  The same argument works at fixed
individually allowed outside paths and with any checks deleted.
\end{proof}

We finish by recording the density cost of prescribing a midpoint.
Under $B_i$ the midpoint density in its local lift is
\begin{equation}\label{eq:bridge-midpoint-density}
                g_i(y)=(2\pi t)^{-3/2}
                         e^{-|y-m_i|^2/(2t)}.
\end{equation}
Replacing it by a uniform midpoint in a unit cell at an absolute
bounded distance from $m_i$ costs at most $CD^{Cx}$ in density.
For two labels with $|m_i-m_k|\le C$ for an absolute constant, the
Gaussian formula gives, throughout either label's pin range,
\begin{equation}\label{eq:midpoint-ratio}
                         \frac{g_i(y)}{g_k(y)}
                                   \le C_xD^{C\sqrt x}.
\end{equation}
Indeed the difference of the two squared distances is at most
$Cb_0$, and $b_0/t=O(\sqrt x\log D)$ with the fixed constants in
\eqref{eq:path-scales}.  Conditioning a pair trial from Assertion~2
on success costs at most one additional factor two.  These density
statements concern a primary path, with the auxiliary bridge kept
conditional on its midpoint if needed.  A pair constrained to share
its midpoint has no density relative to two independent full bridge
laws; no such density will be used.
